\section{问题重述}

\subsection{问题背景}

大语言模型的训练可以类比为培养一名学生：参数量 $N$ 相当于脑容量，训练数据量 $D$ 相当于阅读量，计算量 $C$ 相当于教育经费，数据质量 $Q$ 相当于教材质量，领域配比 $p$ 相当于各科课时分配。标度律刻画损失随规模变化的经验关系，常见形式为 $L(N,D)=E+AN^{-\alpha}+BD^{-\beta}$。然而高质量语料日益稀缺、数据污染加剧，``数据墙''问题凸显；近期研究也表明，参数量小得多的模型可以凭借更高的数据质量在若干任务上胜过巨型模型\cite{asai2026synthesizing}。这说明质量与配比对最终能力的贡献不可忽视。

扩大参数量、增加数据量、提升数据质量和延长上下文长度都需要消耗算力，而算力是有限的。其中上下文长度越大，注意力计算与 KV 缓存等序列相关开销越高，会进一步挤压可用于参数、数据或质量提升的预算。调整领域配比本身不增加总算力开销，却会改变算力在各领域间的分配效果。四者之间存在真实的取舍，这正是本题要研究的核心矛盾。

\subsection{已知条件与数据}

题面提供三类附件，全部位于 \texttt{real\_attachments/} 目录：

\begin{itemize}[leftmargin=2em,itemsep=2pt]
\item \textbf{附件 A}（问题一）：质量信号数据 A1--A3 给出 7 个质量域、22 个质量指标的样本级评分；配方实验数据 A4--A15 给出 17 个训练域（The Pile 子领域）的配比 $p$ 与 13 个验证域的交叉熵损失，其中 A4/A5 为训练集（512 组）、A6--A11 为检验集（1M/60M/1B 三种参数规模）、A12--A15 为外推表；A16 给出两个域体系之间的参考映射。
\item \textbf{附件 B}（问题二）：B1 为 Pythia 训练日志（1176 点），B2 为 Cerebras 半合成轨迹，B3 为插值轨迹，B4/B5 为跨族与文献标度律基准，B6--B8 为含质量维度的半合成实验，B9/B10 为百亿参数以上模型的元数据与估算损失。
\item \textbf{附件 C}（问题三、四）：C1/C2 为 Open LLM Leaderboard 六维评测榜单，C3 为含历史模型的扩展时序，C4 为 Epoch AI 全量模型元数据（含训练算力、数据量、权重开放性），C5/C6 为 Loss--Benchmark 桥接数据，C7 为模型架构元数据（含最大上下文长度），C8 为 1863 个模型的逐任务评测 JSON。
\end{itemize}

\subsection{待求目标}

四个问题构成递进主线，前一问的输出是后一问的输入：

\textbf{问题一：}对 22 个质量指标做预处理与综合评价，给出样本级、语料级或领域级的质量评分 $Q$；定义并消解质量指标之间的冲突；利用 A4--A15 建立 17 域配比 $p$ 与交叉熵损失的定量关系，在检验集上验证并讨论外推稳健性。

\textbf{问题二：}综合附件 A、B 的信息，建立同时包含 $N$、$D$、$Q$、$p$ 的广义标度律；分析各因素的边际效用与弹性，推导质量提升与规模扩大可相互替代的条件。

\textbf{问题三：}在总预算 $C$ 约束下求解 $N$、$D$、$Q$、$p$ 的最优配置，讨论成本函数选择的影响，给出``结构性转移''的数学定义与识别方法，并解析给出上下文长度的临界值 $\ell_{\mathrm{crit}}=6/\eta$。

\textbf{问题四：}分离并量化规模扩张与非规模技术进步各自的贡献占比；在算力增长放缓情景下预测未来 12 或 24 个月开源大语言模型能力的前沿边界，并给出不确定性分析。

\subsection{对题面概念的澄清}

题面中有几处需要在建模前明确口径，本文的处理如下，具体依据见\ref{sec:assume}节：

\textbf{``全量质量信号''的含义。} 题面要求使用 A1 与 A2、A3 的全部记录。经核查，仓库内这四份 \texttt{.jsonl.xz} 文件均为 Git LFS 指针（132--134 字节），且 GitHub 上的 LFS 对象未上传（LFS Batch API 返回 404），实体无法获得。本文按《数据说明》给出的来源（公开数据集 SlimPajama-Meta-rater）同源重建，重建规则与核验结果见\ref{sec:data_integrity}节，并在全文显著标注。

\textbf{``数据质量''的可比性。} 22 个指标的原始量纲与方向均不相同（如教育价值越高越好、广告含量越低越好），题面明确要求统一为``越高越好''。本文在冻结的参考集上估计归一化参数后再做方向补转换，避免对每个域分别做分位数变换后直接比较域级高低。

\textbf{``领域配比''的单纯形约束。} 配比向量满足 $p_i\ge0$、$\sum_i p_i=1$。原始表的行和落在 $0.996\sim1.003$，本文统一做行归一化，且不直接对 $p$ 取对数。

\textbf{``上下文长度''的地位。} 按题面，$\ell$ 是外生情景变量而非内点寻优变量，其可行取值依据 C7 的架构元数据。本文按 C7 中实际出现的 2048/4096/8192/32768/131072 五档取值做情景分析。

\textbf{``开源''的筛选口径。} 题面要求问题四预测开源模型的能力前沿，故须说明入选规则。本文给出严格口径（C4 \texttt{Open model weights?} 字段为 Yes）与宽松口径（榜单 Hub License 非空）两种筛选，分别报告并以宽松口径作为主预测基准。
